\documentclass{report}
\usepackage{amsmath,amssymb}
\setlength{\parindent}{0mm}
\setlength{\parskip}{1em}
\begin{document}
\begin{center}
\rule{6in}{1pt} \
{\large Zichao Di \\
{\bf Applications and Recent Developments of Multilevel Optimization Framework(MG/OPT)}}

5541 Ridgeton Hill Ct Fairfax \\ VA \\ 22032
\\
{\tt zdi@masonlive.gmu.edu}\\
Maria Emelianenko\\
Stephen Nash\end{center}

This talk will present recent progress related to the theory and
applications of the multilevel optimization approach MG/OPT. Inspired by
the traditional multigrid, the intent of MG/OPT is to use calculations on
coarser levels to accelerate the progress of the optimization on the
finest level. There are several advantages of MG/OPT: first of all, the
optimization perspective is broader than systems of equations since it
can handle the constraints in a more natural way. Also, MG/OPT has
stronger guarantees of convergence than traditional multigrid.
Furthermore, in many cases the reduced Hessian of the optimization model
is an elliptic operator making MG/OPT an effective algorithm.
So far, MG/OPT has been successfully applied to many PDE-constrained
optimization problems which have regular geometric structures. However,
in practice, more sophisticated implementations that provide an easier
way to apply MG/OPT to specific problems are still under exploration.
In order to investigate the potential of MG/OPT to solve the problems
with irregular geometric structures, MG/OPT is applied to a particular
type of tessellation problems called centroidal Voronoi
tessellations(CVTs). Significant speedup by using MG/OPT comparing to
other existing techniques has shown.

Furthermore, the constrained version of MG/OPT has also been developed
and it has been applied to successfully solve a variety of other
constrained optimization problems.


\end{document}
